% Figure: flux limiters and three-point stencils at a level transition
% (Chapter 6). Limiter curves sampled from the code's _flux_limiter and the
% upstream neighbours of the code's schedule: generated/momentum_limiter.tex;
% cells: generated/momentum_dual.tex (momentum.jl). Labels hand-written.
\begin{figure}
  \centering
  \begin{subfigure}[b]{0.45\textwidth}
    \centering
    \begin{tikzpicture}[x=1.0cm,y=1.4cm]
      % Second-order TVD region of Sweby (between minmod and superbee) and
      % the TVD region 0 <= Psi <= min(2 Theta, 2).
      \fill[black!6] (0,0) -- (1,2) -- (3,2) -- (3,0) -- cycle;
      \fill[trblue!14] (0,0) -- (0.5,1) -- (1,1) -- (2,2) -- (3,2) -- (3,1) -- (1,1) -- cycle;
      \draw[tr thin vector,black] (0,0) -- (3.2,0) node[anchor=north east,tr small] {$\limratio{g}$};
      \draw[tr thin vector,black] (0,0) -- (0,2.3) node[anchor=north east,tr small] {$\limiter{\mathrm{M}}$};
      \foreach \t in {1,2,3} \draw (\t,0) -- (\t,-0.04) node[below,tr tiny] {\t};
      \foreach \p in {1,2} \draw (0,\p) -- (-0.04,\p) node[left,tr tiny] {\p};
      \draw[black!50,line width=0.4pt] (0,1) -- (3,1);
      \node[tr tiny,anchor=south east,text=black!60] at (3,1) {\code{LaxWendroff}};
      \def\fxlstencil#1#2#3#4#5#6#7#8#9{}%
      \def\fxlcurve#1#2{%
        \ifnum\pdfstrcmp{#1}{minmod}=0 \draw[black,line width=1.1pt] #2;\fi
        \ifnum\pdfstrcmp{#1}{mc}=0 \draw[trred,line width=0.9pt,dash pattern=on 4pt off 1.5pt] #2;\fi
        \ifnum\pdfstrcmp{#1}{vanleer}=0 \draw[trgreen!70!black,line width=0.9pt,dash pattern=on 1pt off 1.2pt] #2;\fi
        \ifnum\pdfstrcmp{#1}{superbee}=0 \draw[trpurple,line width=0.9pt,dash pattern=on 4pt off 1.2pt on 1pt off 1.2pt] #2;\fi
        \ifnum\pdfstrcmp{#1}{fromm}=0 \begin{scope}\clip (0,0) rectangle (3,2.2);
          \draw[trgrey,line width=0.6pt] #2;\end{scope}\fi}%
      % Generated by docs/techref/figures/generate.jl; do not edit.
\fxlcurve{minmod}{(0,0) -- (0.02,0.02) -- (0.04,0.04) -- (0.06,0.06) -- (0.08,0.08) -- (0.1,0.1) -- (0.12,0.12) -- (0.14,0.14) -- (0.16,0.16) -- (0.18,0.18) -- (0.2,0.2) -- (0.22,0.22) -- (0.24,0.24) -- (0.26,0.26) -- (0.28,0.28) -- (0.3,0.3) -- (0.32,0.32) -- (0.34,0.34) -- (0.36,0.36) -- (0.38,0.38) -- (0.4,0.4) -- (0.42,0.42) -- (0.44,0.44) -- (0.46,0.46) -- (0.48,0.48) -- (0.5,0.5) -- (0.52,0.52) -- (0.54,0.54) -- (0.56,0.56) -- (0.58,0.58) -- (0.6,0.6) -- (0.62,0.62) -- (0.64,0.64) -- (0.66,0.66) -- (0.68,0.68) -- (0.7,0.7) -- (0.72,0.72) -- (0.74,0.74) -- (0.76,0.76) -- (0.78,0.78) -- (0.8,0.8) -- (0.82,0.82) -- (0.84,0.84) -- (0.86,0.86) -- (0.88,0.88) -- (0.9,0.9) -- (0.92,0.92) -- (0.94,0.94) -- (0.96,0.96) -- (0.98,0.98) -- (1,1) -- (1.02,1) -- (1.04,1) -- (1.06,1) -- (1.08,1) -- (1.1,1) -- (1.12,1) -- (1.14,1) -- (1.16,1) -- (1.18,1) -- (1.2,1) -- (1.22,1) -- (1.24,1) -- (1.26,1) -- (1.28,1) -- (1.3,1) -- (1.32,1) -- (1.34,1) -- (1.36,1) -- (1.38,1) -- (1.4,1) -- (1.42,1) -- (1.44,1) -- (1.46,1) -- (1.48,1) -- (1.5,1) -- (1.52,1) -- (1.54,1) -- (1.56,1) -- (1.58,1) -- (1.6,1) -- (1.62,1) -- (1.64,1) -- (1.66,1) -- (1.68,1) -- (1.7,1) -- (1.72,1) -- (1.74,1) -- (1.76,1) -- (1.78,1) -- (1.8,1) -- (1.82,1) -- (1.84,1) -- (1.86,1) -- (1.88,1) -- (1.9,1) -- (1.92,1) -- (1.94,1) -- (1.96,1) -- (1.98,1) -- (2,1) -- (2.02,1) -- (2.04,1) -- (2.06,1) -- (2.08,1) -- (2.1,1) -- (2.12,1) -- (2.14,1) -- (2.16,1) -- (2.18,1) -- (2.2,1) -- (2.22,1) -- (2.24,1) -- (2.26,1) -- (2.28,1) -- (2.3,1) -- (2.32,1) -- (2.34,1) -- (2.36,1) -- (2.38,1) -- (2.4,1) -- (2.42,1) -- (2.44,1) -- (2.46,1) -- (2.48,1) -- (2.5,1) -- (2.52,1) -- (2.54,1) -- (2.56,1) -- (2.58,1) -- (2.6,1) -- (2.62,1) -- (2.64,1) -- (2.66,1) -- (2.68,1) -- (2.7,1) -- (2.72,1) -- (2.74,1) -- (2.76,1) -- (2.78,1) -- (2.8,1) -- (2.82,1) -- (2.84,1) -- (2.86,1) -- (2.88,1) -- (2.9,1) -- (2.92,1) -- (2.94,1) -- (2.96,1) -- (2.98,1) -- (3,1)}
\fxlcurve{mc}{(0,0) -- (0.02,0.04) -- (0.04,0.08) -- (0.06,0.12) -- (0.08,0.16) -- (0.1,0.2) -- (0.12,0.24) -- (0.14,0.28) -- (0.16,0.32) -- (0.18,0.36) -- (0.2,0.4) -- (0.22,0.44) -- (0.24,0.48) -- (0.26,0.52) -- (0.28,0.56) -- (0.3,0.6) -- (0.32,0.64) -- (0.34,0.67) -- (0.36,0.68) -- (0.38,0.69) -- (0.4,0.7) -- (0.42,0.71) -- (0.44,0.72) -- (0.46,0.73) -- (0.48,0.74) -- (0.5,0.75) -- (0.52,0.76) -- (0.54,0.77) -- (0.56,0.78) -- (0.58,0.79) -- (0.6,0.8) -- (0.62,0.81) -- (0.64,0.82) -- (0.66,0.83) -- (0.68,0.84) -- (0.7,0.85) -- (0.72,0.86) -- (0.74,0.87) -- (0.76,0.88) -- (0.78,0.89) -- (0.8,0.9) -- (0.82,0.91) -- (0.84,0.92) -- (0.86,0.93) -- (0.88,0.94) -- (0.9,0.95) -- (0.92,0.96) -- (0.94,0.97) -- (0.96,0.98) -- (0.98,0.99) -- (1,1) -- (1.02,1.01) -- (1.04,1.02) -- (1.06,1.03) -- (1.08,1.04) -- (1.1,1.05) -- (1.12,1.06) -- (1.14,1.07) -- (1.16,1.08) -- (1.18,1.09) -- (1.2,1.1) -- (1.22,1.11) -- (1.24,1.12) -- (1.26,1.13) -- (1.28,1.14) -- (1.3,1.15) -- (1.32,1.16) -- (1.34,1.17) -- (1.36,1.18) -- (1.38,1.19) -- (1.4,1.2) -- (1.42,1.21) -- (1.44,1.22) -- (1.46,1.23) -- (1.48,1.24) -- (1.5,1.25) -- (1.52,1.26) -- (1.54,1.27) -- (1.56,1.28) -- (1.58,1.29) -- (1.6,1.3) -- (1.62,1.31) -- (1.64,1.32) -- (1.66,1.33) -- (1.68,1.34) -- (1.7,1.35) -- (1.72,1.36) -- (1.74,1.37) -- (1.76,1.38) -- (1.78,1.39) -- (1.8,1.4) -- (1.82,1.41) -- (1.84,1.42) -- (1.86,1.43) -- (1.88,1.44) -- (1.9,1.45) -- (1.92,1.46) -- (1.94,1.47) -- (1.96,1.48) -- (1.98,1.49) -- (2,1.5) -- (2.02,1.51) -- (2.04,1.52) -- (2.06,1.53) -- (2.08,1.54) -- (2.1,1.55) -- (2.12,1.56) -- (2.14,1.57) -- (2.16,1.58) -- (2.18,1.59) -- (2.2,1.6) -- (2.22,1.61) -- (2.24,1.62) -- (2.26,1.63) -- (2.28,1.64) -- (2.3,1.65) -- (2.32,1.66) -- (2.34,1.67) -- (2.36,1.68) -- (2.38,1.69) -- (2.4,1.7) -- (2.42,1.71) -- (2.44,1.72) -- (2.46,1.73) -- (2.48,1.74) -- (2.5,1.75) -- (2.52,1.76) -- (2.54,1.77) -- (2.56,1.78) -- (2.58,1.79) -- (2.6,1.8) -- (2.62,1.81) -- (2.64,1.82) -- (2.66,1.83) -- (2.68,1.84) -- (2.7,1.85) -- (2.72,1.86) -- (2.74,1.87) -- (2.76,1.88) -- (2.78,1.89) -- (2.8,1.9) -- (2.82,1.91) -- (2.84,1.92) -- (2.86,1.93) -- (2.88,1.94) -- (2.9,1.95) -- (2.92,1.96) -- (2.94,1.97) -- (2.96,1.98) -- (2.98,1.99) -- (3,2)}
\fxlcurve{vanleer}{(0,0) -- (0.02,0.039216) -- (0.04,0.076923) -- (0.06,0.113208) -- (0.08,0.148148) -- (0.1,0.181818) -- (0.12,0.214286) -- (0.14,0.245614) -- (0.16,0.275862) -- (0.18,0.305085) -- (0.2,0.333333) -- (0.22,0.360656) -- (0.24,0.387097) -- (0.26,0.412698) -- (0.28,0.4375) -- (0.3,0.461538) -- (0.32,0.484848) -- (0.34,0.507463) -- (0.36,0.529412) -- (0.38,0.550725) -- (0.4,0.571429) -- (0.42,0.591549) -- (0.44,0.611111) -- (0.46,0.630137) -- (0.48,0.648649) -- (0.5,0.666667) -- (0.52,0.684211) -- (0.54,0.701299) -- (0.56,0.717949) -- (0.58,0.734177) -- (0.6,0.75) -- (0.62,0.765432) -- (0.64,0.780488) -- (0.66,0.795181) -- (0.68,0.809524) -- (0.7,0.823529) -- (0.72,0.837209) -- (0.74,0.850575) -- (0.76,0.863636) -- (0.78,0.876404) -- (0.8,0.888889) -- (0.82,0.901099) -- (0.84,0.913043) -- (0.86,0.924731) -- (0.88,0.93617) -- (0.9,0.947368) -- (0.92,0.958333) -- (0.94,0.969072) -- (0.96,0.979592) -- (0.98,0.989899) -- (1,1) -- (1.02,1.009901) -- (1.04,1.019608) -- (1.06,1.029126) -- (1.08,1.038462) -- (1.1,1.047619) -- (1.12,1.056604) -- (1.14,1.065421) -- (1.16,1.074074) -- (1.18,1.082569) -- (1.2,1.090909) -- (1.22,1.099099) -- (1.24,1.107143) -- (1.26,1.115044) -- (1.28,1.122807) -- (1.3,1.130435) -- (1.32,1.137931) -- (1.34,1.145299) -- (1.36,1.152542) -- (1.38,1.159664) -- (1.4,1.166667) -- (1.42,1.173554) -- (1.44,1.180328) -- (1.46,1.186992) -- (1.48,1.193548) -- (1.5,1.2) -- (1.52,1.206349) -- (1.54,1.212598) -- (1.56,1.21875) -- (1.58,1.224806) -- (1.6,1.230769) -- (1.62,1.236641) -- (1.64,1.242424) -- (1.66,1.24812) -- (1.68,1.253731) -- (1.7,1.259259) -- (1.72,1.264706) -- (1.74,1.270073) -- (1.76,1.275362) -- (1.78,1.280576) -- (1.8,1.285714) -- (1.82,1.29078) -- (1.84,1.295775) -- (1.86,1.300699) -- (1.88,1.305556) -- (1.9,1.310345) -- (1.92,1.315068) -- (1.94,1.319728) -- (1.96,1.324324) -- (1.98,1.328859) -- (2,1.333333) -- (2.02,1.337748) -- (2.04,1.342105) -- (2.06,1.346405) -- (2.08,1.350649) -- (2.1,1.354839) -- (2.12,1.358974) -- (2.14,1.363057) -- (2.16,1.367089) -- (2.18,1.371069) -- (2.2,1.375) -- (2.22,1.378882) -- (2.24,1.382716) -- (2.26,1.386503) -- (2.28,1.390244) -- (2.3,1.393939) -- (2.32,1.39759) -- (2.34,1.401198) -- (2.36,1.404762) -- (2.38,1.408284) -- (2.4,1.411765) -- (2.42,1.415205) -- (2.44,1.418605) -- (2.46,1.421965) -- (2.48,1.425287) -- (2.5,1.428571) -- (2.52,1.431818) -- (2.54,1.435028) -- (2.56,1.438202) -- (2.58,1.441341) -- (2.6,1.444444) -- (2.62,1.447514) -- (2.64,1.450549) -- (2.66,1.453552) -- (2.68,1.456522) -- (2.7,1.459459) -- (2.72,1.462366) -- (2.74,1.465241) -- (2.76,1.468085) -- (2.78,1.470899) -- (2.8,1.473684) -- (2.82,1.47644) -- (2.84,1.479167) -- (2.86,1.481865) -- (2.88,1.484536) -- (2.9,1.487179) -- (2.92,1.489796) -- (2.94,1.492386) -- (2.96,1.494949) -- (2.98,1.497487) -- (3,1.5)}
\fxlcurve{superbee}{(0,0) -- (0.02,0.04) -- (0.04,0.08) -- (0.06,0.12) -- (0.08,0.16) -- (0.1,0.2) -- (0.12,0.24) -- (0.14,0.28) -- (0.16,0.32) -- (0.18,0.36) -- (0.2,0.4) -- (0.22,0.44) -- (0.24,0.48) -- (0.26,0.52) -- (0.28,0.56) -- (0.3,0.6) -- (0.32,0.64) -- (0.34,0.68) -- (0.36,0.72) -- (0.38,0.76) -- (0.4,0.8) -- (0.42,0.84) -- (0.44,0.88) -- (0.46,0.92) -- (0.48,0.96) -- (0.5,1) -- (0.52,1) -- (0.54,1) -- (0.56,1) -- (0.58,1) -- (0.6,1) -- (0.62,1) -- (0.64,1) -- (0.66,1) -- (0.68,1) -- (0.7,1) -- (0.72,1) -- (0.74,1) -- (0.76,1) -- (0.78,1) -- (0.8,1) -- (0.82,1) -- (0.84,1) -- (0.86,1) -- (0.88,1) -- (0.9,1) -- (0.92,1) -- (0.94,1) -- (0.96,1) -- (0.98,1) -- (1,1) -- (1.02,1.02) -- (1.04,1.04) -- (1.06,1.06) -- (1.08,1.08) -- (1.1,1.1) -- (1.12,1.12) -- (1.14,1.14) -- (1.16,1.16) -- (1.18,1.18) -- (1.2,1.2) -- (1.22,1.22) -- (1.24,1.24) -- (1.26,1.26) -- (1.28,1.28) -- (1.3,1.3) -- (1.32,1.32) -- (1.34,1.34) -- (1.36,1.36) -- (1.38,1.38) -- (1.4,1.4) -- (1.42,1.42) -- (1.44,1.44) -- (1.46,1.46) -- (1.48,1.48) -- (1.5,1.5) -- (1.52,1.52) -- (1.54,1.54) -- (1.56,1.56) -- (1.58,1.58) -- (1.6,1.6) -- (1.62,1.62) -- (1.64,1.64) -- (1.66,1.66) -- (1.68,1.68) -- (1.7,1.7) -- (1.72,1.72) -- (1.74,1.74) -- (1.76,1.76) -- (1.78,1.78) -- (1.8,1.8) -- (1.82,1.82) -- (1.84,1.84) -- (1.86,1.86) -- (1.88,1.88) -- (1.9,1.9) -- (1.92,1.92) -- (1.94,1.94) -- (1.96,1.96) -- (1.98,1.98) -- (2,2) -- (2.02,2) -- (2.04,2) -- (2.06,2) -- (2.08,2) -- (2.1,2) -- (2.12,2) -- (2.14,2) -- (2.16,2) -- (2.18,2) -- (2.2,2) -- (2.22,2) -- (2.24,2) -- (2.26,2) -- (2.28,2) -- (2.3,2) -- (2.32,2) -- (2.34,2) -- (2.36,2) -- (2.38,2) -- (2.4,2) -- (2.42,2) -- (2.44,2) -- (2.46,2) -- (2.48,2) -- (2.5,2) -- (2.52,2) -- (2.54,2) -- (2.56,2) -- (2.58,2) -- (2.6,2) -- (2.62,2) -- (2.64,2) -- (2.66,2) -- (2.68,2) -- (2.7,2) -- (2.72,2) -- (2.74,2) -- (2.76,2) -- (2.78,2) -- (2.8,2) -- (2.82,2) -- (2.84,2) -- (2.86,2) -- (2.88,2) -- (2.9,2) -- (2.92,2) -- (2.94,2) -- (2.96,2) -- (2.98,2) -- (3,2)}
\fxlcurve{fromm}{(0,0.5) -- (0.02,0.51) -- (0.04,0.52) -- (0.06,0.53) -- (0.08,0.54) -- (0.1,0.55) -- (0.12,0.56) -- (0.14,0.57) -- (0.16,0.58) -- (0.18,0.59) -- (0.2,0.6) -- (0.22,0.61) -- (0.24,0.62) -- (0.26,0.63) -- (0.28,0.64) -- (0.3,0.65) -- (0.32,0.66) -- (0.34,0.67) -- (0.36,0.68) -- (0.38,0.69) -- (0.4,0.7) -- (0.42,0.71) -- (0.44,0.72) -- (0.46,0.73) -- (0.48,0.74) -- (0.5,0.75) -- (0.52,0.76) -- (0.54,0.77) -- (0.56,0.78) -- (0.58,0.79) -- (0.6,0.8) -- (0.62,0.81) -- (0.64,0.82) -- (0.66,0.83) -- (0.68,0.84) -- (0.7,0.85) -- (0.72,0.86) -- (0.74,0.87) -- (0.76,0.88) -- (0.78,0.89) -- (0.8,0.9) -- (0.82,0.91) -- (0.84,0.92) -- (0.86,0.93) -- (0.88,0.94) -- (0.9,0.95) -- (0.92,0.96) -- (0.94,0.97) -- (0.96,0.98) -- (0.98,0.99) -- (1,1) -- (1.02,1.01) -- (1.04,1.02) -- (1.06,1.03) -- (1.08,1.04) -- (1.1,1.05) -- (1.12,1.06) -- (1.14,1.07) -- (1.16,1.08) -- (1.18,1.09) -- (1.2,1.1) -- (1.22,1.11) -- (1.24,1.12) -- (1.26,1.13) -- (1.28,1.14) -- (1.3,1.15) -- (1.32,1.16) -- (1.34,1.17) -- (1.36,1.18) -- (1.38,1.19) -- (1.4,1.2) -- (1.42,1.21) -- (1.44,1.22) -- (1.46,1.23) -- (1.48,1.24) -- (1.5,1.25) -- (1.52,1.26) -- (1.54,1.27) -- (1.56,1.28) -- (1.58,1.29) -- (1.6,1.3) -- (1.62,1.31) -- (1.64,1.32) -- (1.66,1.33) -- (1.68,1.34) -- (1.7,1.35) -- (1.72,1.36) -- (1.74,1.37) -- (1.76,1.38) -- (1.78,1.39) -- (1.8,1.4) -- (1.82,1.41) -- (1.84,1.42) -- (1.86,1.43) -- (1.88,1.44) -- (1.9,1.45) -- (1.92,1.46) -- (1.94,1.47) -- (1.96,1.48) -- (1.98,1.49) -- (2,1.5) -- (2.02,1.51) -- (2.04,1.52) -- (2.06,1.53) -- (2.08,1.54) -- (2.1,1.55) -- (2.12,1.56) -- (2.14,1.57) -- (2.16,1.58) -- (2.18,1.59) -- (2.2,1.6) -- (2.22,1.61) -- (2.24,1.62) -- (2.26,1.63) -- (2.28,1.64) -- (2.3,1.65) -- (2.32,1.66) -- (2.34,1.67) -- (2.36,1.68) -- (2.38,1.69) -- (2.4,1.7) -- (2.42,1.71) -- (2.44,1.72) -- (2.46,1.73) -- (2.48,1.74) -- (2.5,1.75) -- (2.52,1.76) -- (2.54,1.77) -- (2.56,1.78) -- (2.58,1.79) -- (2.6,1.8) -- (2.62,1.81) -- (2.64,1.82) -- (2.66,1.83) -- (2.68,1.84) -- (2.7,1.85) -- (2.72,1.86) -- (2.74,1.87) -- (2.76,1.88) -- (2.78,1.89) -- (2.8,1.9) -- (2.82,1.91) -- (2.84,1.92) -- (2.86,1.93) -- (2.88,1.94) -- (2.9,1.95) -- (2.92,1.96) -- (2.94,1.97) -- (2.96,1.98) -- (2.98,1.99) -- (3,2)}
\fxlstencil{4}{5.25}{3}{5.5}{4.5}{5.25}{5}{5}{2}
\fxlstencil{4}{5.25}{4.5}{5.25}{3}{5.5}{3.5}{5.25}{\tfrac{1}{2}}

      \node[tr tiny,anchor=west] at (3.03,1.0) {\code{Minmod}};
      \node[tr tiny,anchor=west,text=trred] at (3.03,2.08) {\code{MonotonizedCentral}};
      \node[tr tiny,anchor=west,text=trgreen!50!black] at (3.03,1.5) {\code{VanLeer}};
      \node[tr tiny,anchor=west,text=trpurple] at (3.03,1.92) {\code{Superbee}};
      \node[tr tiny,anchor=south,text=trgrey,rotate=26.6] at (2.2,1.66) {\code{Fromm}};
    \end{tikzpicture}
    \caption{Limiters $\limiter{\mathrm{M}}(\limratio{g})$ (shaded: the TVD
      region, blue: its second-order part).}
    \label{fig:momentum:limiter:sweby}
  \end{subfigure}%
  \hfill
  \begin{subfigure}[b]{0.53\textwidth}
    \centering
    \begin{tikzpicture}[x=2.6cm,y=2.6cm]
      \def\fxmnone{%
        \def\fxmcell##1##2##3##4##5##6##7{}\def\fxmsolid##1##2{}%
        \def\fxmwall##1##2##3##4##5{}\def\fxmdual##1##2##3##4##5##6{}%
        \def\fxmopen##1##2{}\def\fxmface##1##2##3##4##5##6##7{}%
        \def\fxmlink##1##2##3##4##5##6##7##8##9{}%
        \def\fxmdualface##1##2##3##4##5##6##7{}\def\fxmhanging##1##2##3{}}%
      \begin{scope}
        \clip (2.6,4.75) rectangle (5.0,6.25);
        \fxmnone
        \def\fxmcell#1#2#3#4#5#6#7{\ifnum#1=2
          \fill[tr level #6] (#2,#3) rectangle (#4,#5);
          \draw[tr cell] (#2,#3) rectangle (#4,#5);\fi}%
        % Generated by docs/techref/figures/generate.jl; do not edit.
\fxmcell{1}{4}{4}{5}{5}{0}{0}
\fxmcell{1}{5}{4}{6}{5}{0}{0}
\fxmcell{1}{6}{4}{7}{5}{0}{0}
\fxmcell{1}{7}{4}{8}{5}{0}{0}
\fxmcell{1}{4}{5}{5}{6}{0}{0}
\fxmcell{1}{5}{5}{6}{6}{0}{0}
\fxmcell{1}{6}{5}{7}{6}{0}{0}
\fxmcell{1}{7}{5}{8}{6}{0}{0}
\fxmcell{1}{4}{6}{5}{7}{0}{0}
\fxmcell{1}{5}{6}{6}{7}{0}{0}
\fxmcell{1}{6}{6}{7}{7}{0}{0}
\fxmcell{1}{7}{6}{8}{7}{0}{0}
\fxmdual{1}{4.5}{4}{5.5}{5}{0}
\fxmdual{1}{4.5}{5}{5.5}{6}{0}
\fxmdual{1}{4.5}{6}{5.5}{7}{0}
\fxmdual{1}{5.5}{4}{6.5}{5}{0}
\fxmdual{1}{5.5}{5}{6.5}{6}{1}
\fxmdual{1}{5.5}{6}{6.5}{7}{0}
\fxmdual{1}{6.5}{4}{7.5}{5}{0}
\fxmdual{1}{6.5}{5}{7.5}{6}{0}
\fxmdual{1}{6.5}{6}{7.5}{7}{0}
\fxmface{1}{1}{6}{5.5}{1}{f}{1.00}
\fxmface{1}{2}{6.5}{6}{1}{k'}{1.00}
\fxmface{1}{2}{5.5}{5}{1}{l}{1.00}
\fxmface{1}{2}{6.5}{5}{1}{l'}{1.00}
\fxmface{1}{1}{7}{5.5}{1}{g'}{1.00}
\fxmface{1}{1}{5}{5.5}{1}{g}{1.00}
\fxmface{1}{2}{5.5}{6}{1}{k}{1.00}
\fxmdualface{1}{5.5}{5}{5.5}{6}{\tfrac{1}{2}\,f+\tfrac{1}{2}\,g}{l}
\fxmdualface{1}{6.5}{5}{6.5}{6}{\tfrac{1}{2}\,f+\tfrac{1}{2}\,g'}{r}
\fxmdualface{1}{5.5}{5}{6.5}{5}{\tfrac{1}{2}\,l+\tfrac{1}{2}\,l'}{b}
\fxmdualface{1}{5.5}{6}{6.5}{6}{\tfrac{1}{2}\,k+\tfrac{1}{2}\,k'}{t}
\fxmcell{2}{2}{3}{3}{4}{0}{0}
\fxmcell{2}{3}{3}{4}{4}{0}{0}
\fxmcell{2}{4}{3}{5}{4}{0}{0}
\fxmcell{2}{5}{3}{6}{4}{0}{0}
\fxmcell{2}{2}{4}{3}{5}{0}{0}
\fxmcell{2}{3}{4}{4}{5}{0}{0}
\fxmcell{2}{2}{5}{3}{6}{0}{0}
\fxmcell{2}{3}{5}{4}{6}{0}{0}
\fxmcell{2}{2}{6}{3}{7}{0}{0}
\fxmcell{2}{3}{6}{4}{7}{0}{0}
\fxmcell{2}{2}{7}{3}{8}{0}{0}
\fxmcell{2}{3}{7}{4}{8}{0}{0}
\fxmcell{2}{4}{4}{4.5}{4.5}{1}{0}
\fxmcell{2}{4.5}{4}{5}{4.5}{1}{0}
\fxmcell{2}{5}{4}{5.5}{4.5}{1}{0}
\fxmcell{2}{4}{4.5}{4.5}{5}{1}{0}
\fxmcell{2}{4.5}{4.5}{5}{5}{1}{0}
\fxmcell{2}{5}{4.5}{5.5}{5}{1}{0}
\fxmcell{2}{4}{5}{4.5}{5.5}{1}{0}
\fxmcell{2}{4.5}{5}{5}{5.5}{1}{0}
\fxmcell{2}{5}{5}{5.5}{5.5}{1}{0}
\fxmcell{2}{4}{5.5}{4.5}{6}{1}{0}
\fxmcell{2}{4.5}{5.5}{5}{6}{1}{0}
\fxmcell{2}{5}{5.5}{5.5}{6}{1}{0}
\fxmcell{2}{4}{6}{4.5}{6.5}{1}{0}
\fxmcell{2}{4.5}{6}{5}{6.5}{1}{0}
\fxmcell{2}{5}{6}{5.5}{6.5}{1}{0}
\fxmcell{2}{4}{6.5}{4.5}{7}{1}{0}
\fxmcell{2}{4.5}{6.5}{5}{7}{1}{0}
\fxmcell{2}{5}{6.5}{5.5}{7}{1}{0}
\fxmcell{2}{4}{7}{4.5}{7.5}{1}{0}
\fxmcell{2}{4.5}{7}{5}{7.5}{1}{0}
\fxmcell{2}{5}{7}{5.5}{7.5}{1}{0}
\fxmdual{2}{1.5}{3}{2.5}{4}{0}
\fxmdual{2}{1.5}{4}{2.5}{5}{0}
\fxmdual{2}{1.5}{5}{2.5}{6}{0}
\fxmdual{2}{1.5}{6}{2.5}{7}{0}
\fxmdual{2}{1.5}{7}{2.5}{8}{0}
\fxmdual{2}{2.5}{3}{3.5}{4}{0}
\fxmdual{2}{2.5}{4}{3.5}{5}{0}
\fxmdual{2}{2.5}{5}{3.5}{6}{0}
\fxmdual{2}{2.5}{6}{3.5}{7}{0}
\fxmdual{2}{2.5}{7}{3.5}{8}{0}
\fxmdual{2}{3.5}{3}{4.5}{4}{0}
\fxmdual{2}{3.5}{4}{4.25}{4.5}{0}
\fxmdual{2}{3.5}{4.5}{4.25}{5}{0}
\fxmdual{2}{3.5}{5}{4.25}{5.5}{1}
\fxmdual{2}{3.5}{5.5}{4.25}{6}{0}
\fxmdual{2}{3.5}{6}{4.25}{6.5}{0}
\fxmdual{2}{3.5}{6.5}{4.25}{7}{0}
\fxmdual{2}{3.5}{7}{4.25}{7.5}{0}
\fxmdual{2}{4.25}{4}{4.75}{4.5}{0}
\fxmdual{2}{4.25}{4.5}{4.75}{5}{0}
\fxmdual{2}{4.25}{5}{4.75}{5.5}{0}
\fxmdual{2}{4.25}{5.5}{4.75}{6}{0}
\fxmdual{2}{4.25}{6}{4.75}{6.5}{0}
\fxmdual{2}{4.25}{6.5}{4.75}{7}{0}
\fxmdual{2}{4.25}{7}{4.75}{7.5}{0}
\fxmdual{2}{4.5}{3}{5.5}{4}{0}
\fxmdual{2}{4.75}{4}{5.25}{4.5}{0}
\fxmdual{2}{4.75}{4.5}{5.25}{5}{0}
\fxmdual{2}{4.75}{5}{5.25}{5.5}{0}
\fxmdual{2}{4.75}{5.5}{5.25}{6}{0}
\fxmdual{2}{4.75}{6}{5.25}{6.5}{0}
\fxmdual{2}{4.75}{6.5}{5.25}{7}{0}
\fxmdual{2}{4.75}{7}{5.25}{7.5}{0}
\fxmface{2}{1}{4}{5.75}{0.5}{f_2}{1.00}
\fxmface{2}{1}{4}{5.25}{0.5}{f_1}{1.00}
\fxmface{2}{1}{4.5}{5.25}{0.5}{h_1}{1.00}
\fxmface{2}{2}{3.5}{6}{1}{\bar{k}^{+}}{1.00}
\fxmface{2}{2}{4.25}{5.5}{0.5}{k}{1.00}
\fxmface{2}{1}{3}{5.5}{1}{g}{1.00}
\fxmface{2}{2}{3.5}{5}{1}{\bar{k}^{-}}{1.00}
\fxmface{2}{2}{4.25}{5}{0.5}{k_1}{1.00}
\fxmlink{2}{4}{5.25}{3}{5.5}{1}{1}{\tfrac{1}{2}\,f_1+\tfrac{1}{4}\,g}{0}
\fxmlink{2}{4}{5.25}{4}{4.25}{2}{1}{\tfrac{1}{8}\,\bar{k}^{-}}{-1}
\fxmlink{2}{4}{5.25}{4}{4.75}{2}{1}{\tfrac{1}{8}\,\bar{k}^{-}+\tfrac{1}{2}\,k_1}{-1}
\fxmlink{2}{4}{5.25}{4}{5.75}{1}{1}{-\tfrac{1}{4}\,f_1+\tfrac{1}{4}\,f_2}{1}
\fxmlink{2}{4}{5.25}{4}{5.75}{2}{1}{\tfrac{1}{2}\,k}{-1}
\fxmlink{2}{4}{5.25}{4}{6.25}{2}{1}{\tfrac{1}{8}\,\bar{k}^{+}}{-1}
\fxmlink{2}{4}{5.25}{4}{6.75}{2}{1}{\tfrac{1}{8}\,\bar{k}^{+}}{-1}
\fxmlink{2}{4}{5.25}{4.5}{5.25}{1}{1}{\tfrac{1}{2}\,f_1+\tfrac{1}{2}\,h_1}{0}
\fxmhanging{2}{4}{4.5}
\fxmhanging{2}{4}{4.5}
\fxmhanging{2}{4}{5.5}
\fxmhanging{2}{4}{5.5}
\fxmhanging{2}{4}{6.5}
\fxmhanging{2}{4}{6.5}
\fxmcell{3}{6}{4.5}{6.5}{5}{1}{0}
\fxmcell{3}{6.5}{4.5}{7}{5}{1}{0}
\fxmcell{3}{7}{4.5}{7.5}{5}{1}{1}
\fxmsolid{3}{(7.5,4.894926) -- (7.5,5) -- (7.244029,5) -- cycle}
\fxmwall{3}{7.5}{4.894926}{7.244029}{5}
\fxmcell{3}{7.5}{4.5}{8}{5}{1}{1}
\fxmsolid{3}{(8,4.783878) -- (8,5) -- (7.5,5) -- (7.5,4.857327) -- cycle}
\fxmwall{3}{8}{4.783878}{7.5}{4.857327}
\fxmcell{3}{6}{5}{6.5}{5.5}{1}{0}
\fxmcell{3}{6.5}{5}{7}{5.5}{1}{1}
\fxmsolid{3}{(7,5.165022) -- (7,5.5) -- (6.774767,5.5) -- cycle}
\fxmwall{3}{7}{5.165022}{6.774767}{5.5}
\fxmcell{3}{7}{5}{7.5}{5.5}{1}{1}
\fxmsolid{3}{(7.22103,5) -- (7.5,5) -- (7.5,5.5) -- (7,5.5) -- (7,5.238627) -- cycle}
\fxmwall{3}{7.22103}{5}{7}{5.238627}
\fxmcell{3}{7.5}{5}{8}{5.5}{1}{2}
\fxmsolid{3}{(7.5,5) -- (8,5) -- (8,5.5) -- (7.5,5.5) -- cycle}
\fxmcell{3}{6}{5.5}{6.5}{6}{1}{0}
\fxmcell{3}{6.5}{5.5}{7}{6}{1}{1}
\fxmsolid{3}{(6.753218,5.5) -- (7,5.5) -- (7,6) -- (6.528991,6) -- cycle}
\fxmwall{3}{6.753218}{5.5}{6.528991}{6}
\fxmcell{3}{7}{5.5}{7.5}{6}{1}{2}
\fxmsolid{3}{(7,5.5) -- (7.5,5.5) -- (7.5,6) -- (7,6) -- cycle}
\fxmcell{3}{7.5}{5.5}{8}{6}{1}{2}
\fxmsolid{3}{(7.5,5.5) -- (8,5.5) -- (8,6) -- (7.5,6) -- cycle}
\fxmopen{3}{(6.5,5) -- (7,5) -- (7,5.165022) -- (6.774767,5.5) -- (6.5,5.5) -- cycle}
\fxmopen{3}{(7,5) -- (7.22103,5) -- (7,5.238627) -- cycle}
\fxmdual{3}{6.75}{5}{7.25}{5.5}{1}
\fxmface{3}{2}{7.25}{5}{0.5}{l'}{0.47}
\fxmface{3}{2}{6.75}{5.5}{0.5}{k}{0.53}
\fxmface{3}{2}{7.25}{5.5}{0.5}{k'}{0.00}
\fxmface{3}{1}{7.5}{5.25}{0.5}{g'}{0.00}
\fxmface{3}{1}{6.5}{5.25}{0.5}{g}{1.00}
\fxmface{3}{1}{7}{5.25}{0.5}{f}{0.40}
\fxmface{3}{2}{6.75}{5}{0.5}{l}{1.00}
\fxmdualface{3}{6.75}{5}{6.75}{5.5}{\tfrac12(\aperture{f}\,f+\aperture{g}\,g)}{l}
\fxmdualface{3}{6.75}{5}{7.25}{5}{\tfrac12(\aperture{l}\,l+\aperture{l'}\,l')}{b}
\fxmdualface{3}{6.75}{5.5}{7.25}{5.5}{\tfrac12(\aperture{k}\,k+\aperture{k'}\,k')}{t}
\fxmdualface{3}{7.25}{5}{7.25}{5.5}{\tfrac12(\aperture{f}\,f+\aperture{g'}\,g')}{r}

      \end{scope}
      \def\fxlcurve#1#2{}%
      % The stencil of the flux towards +e_1 through the dual face (f_1, h_1).
      \def\fxlstencil#1#2#3#4#5#6#7#8#9{%
        \ifdim#3pt>#1pt
          \draw[tr link] (#1,#2) -- (#3,#4);
          \draw[black!60,dash pattern=on 1.5pt off 1.2pt] (#1,#2) -- (#7,#8);
          \node[tr point,fill=trred] (u) at (#1,#2) {};
          \node[tr point,fill=trblue] (d) at (#3,#4) {};
          \node[tr point,fill=trgreen!70!black] (uu) at (#5,#6) {};
          \node[tr vertex,minimum size=3.4pt] (r) at (#7,#8) {};
          \node[tr small,anchor=south west,text=trred] at (u.north east) {$\stencildof{u}{g}$};
          \node[tr small,anchor=south,text=trblue] at (d.north) {$\stencildof{d}{g}$};
          \node[tr small,anchor=south,text=trgreen!50!black] at (uu.north) {$\stencildof{uu}{g}$};
          \draw[decorate,decoration={brace,amplitude=3pt,mirror}] (#1,#2-0.12) -- (#3,#4-0.12)
            node[midway,below=2pt,tr small] {$\axdist{d}{g}$};
          \draw[decorate,decoration={brace,amplitude=3pt,mirror}] (#5,#2-0.12) -- (#1,#2-0.12)
            node[midway,below=2pt,tr small] {$\axdist{uu}{g}$};
          \node[tr tiny,anchor=west,fill=white,fill opacity=0.8,text opacity=1] at (4.56,5.02)
            {$\axdist{d}{g}/\axdist{uu}{g}=#9$};
        \fi}%
      % Generated by docs/techref/figures/generate.jl; do not edit.
\fxlcurve{minmod}{(0,0) -- (0.02,0.02) -- (0.04,0.04) -- (0.06,0.06) -- (0.08,0.08) -- (0.1,0.1) -- (0.12,0.12) -- (0.14,0.14) -- (0.16,0.16) -- (0.18,0.18) -- (0.2,0.2) -- (0.22,0.22) -- (0.24,0.24) -- (0.26,0.26) -- (0.28,0.28) -- (0.3,0.3) -- (0.32,0.32) -- (0.34,0.34) -- (0.36,0.36) -- (0.38,0.38) -- (0.4,0.4) -- (0.42,0.42) -- (0.44,0.44) -- (0.46,0.46) -- (0.48,0.48) -- (0.5,0.5) -- (0.52,0.52) -- (0.54,0.54) -- (0.56,0.56) -- (0.58,0.58) -- (0.6,0.6) -- (0.62,0.62) -- (0.64,0.64) -- (0.66,0.66) -- (0.68,0.68) -- (0.7,0.7) -- (0.72,0.72) -- (0.74,0.74) -- (0.76,0.76) -- (0.78,0.78) -- (0.8,0.8) -- (0.82,0.82) -- (0.84,0.84) -- (0.86,0.86) -- (0.88,0.88) -- (0.9,0.9) -- (0.92,0.92) -- (0.94,0.94) -- (0.96,0.96) -- (0.98,0.98) -- (1,1) -- (1.02,1) -- (1.04,1) -- (1.06,1) -- (1.08,1) -- (1.1,1) -- (1.12,1) -- (1.14,1) -- (1.16,1) -- (1.18,1) -- (1.2,1) -- (1.22,1) -- (1.24,1) -- (1.26,1) -- (1.28,1) -- (1.3,1) -- (1.32,1) -- (1.34,1) -- (1.36,1) -- (1.38,1) -- (1.4,1) -- (1.42,1) -- (1.44,1) -- (1.46,1) -- (1.48,1) -- (1.5,1) -- (1.52,1) -- (1.54,1) -- (1.56,1) -- (1.58,1) -- (1.6,1) -- (1.62,1) -- (1.64,1) -- (1.66,1) -- (1.68,1) -- (1.7,1) -- (1.72,1) -- (1.74,1) -- (1.76,1) -- (1.78,1) -- (1.8,1) -- (1.82,1) -- (1.84,1) -- (1.86,1) -- (1.88,1) -- (1.9,1) -- (1.92,1) -- (1.94,1) -- (1.96,1) -- (1.98,1) -- (2,1) -- (2.02,1) -- (2.04,1) -- (2.06,1) -- (2.08,1) -- (2.1,1) -- (2.12,1) -- (2.14,1) -- (2.16,1) -- (2.18,1) -- (2.2,1) -- (2.22,1) -- (2.24,1) -- (2.26,1) -- (2.28,1) -- (2.3,1) -- (2.32,1) -- (2.34,1) -- (2.36,1) -- (2.38,1) -- (2.4,1) -- (2.42,1) -- (2.44,1) -- (2.46,1) -- (2.48,1) -- (2.5,1) -- (2.52,1) -- (2.54,1) -- (2.56,1) -- (2.58,1) -- (2.6,1) -- (2.62,1) -- (2.64,1) -- (2.66,1) -- (2.68,1) -- (2.7,1) -- (2.72,1) -- (2.74,1) -- (2.76,1) -- (2.78,1) -- (2.8,1) -- (2.82,1) -- (2.84,1) -- (2.86,1) -- (2.88,1) -- (2.9,1) -- (2.92,1) -- (2.94,1) -- (2.96,1) -- (2.98,1) -- (3,1)}
\fxlcurve{mc}{(0,0) -- (0.02,0.04) -- (0.04,0.08) -- (0.06,0.12) -- (0.08,0.16) -- (0.1,0.2) -- (0.12,0.24) -- (0.14,0.28) -- (0.16,0.32) -- (0.18,0.36) -- (0.2,0.4) -- (0.22,0.44) -- (0.24,0.48) -- (0.26,0.52) -- (0.28,0.56) -- (0.3,0.6) -- (0.32,0.64) -- (0.34,0.67) -- (0.36,0.68) -- (0.38,0.69) -- (0.4,0.7) -- (0.42,0.71) -- (0.44,0.72) -- (0.46,0.73) -- (0.48,0.74) -- (0.5,0.75) -- (0.52,0.76) -- (0.54,0.77) -- (0.56,0.78) -- (0.58,0.79) -- (0.6,0.8) -- (0.62,0.81) -- (0.64,0.82) -- (0.66,0.83) -- (0.68,0.84) -- (0.7,0.85) -- (0.72,0.86) -- (0.74,0.87) -- (0.76,0.88) -- (0.78,0.89) -- (0.8,0.9) -- (0.82,0.91) -- (0.84,0.92) -- (0.86,0.93) -- (0.88,0.94) -- (0.9,0.95) -- (0.92,0.96) -- (0.94,0.97) -- (0.96,0.98) -- (0.98,0.99) -- (1,1) -- (1.02,1.01) -- (1.04,1.02) -- (1.06,1.03) -- (1.08,1.04) -- (1.1,1.05) -- (1.12,1.06) -- (1.14,1.07) -- (1.16,1.08) -- (1.18,1.09) -- (1.2,1.1) -- (1.22,1.11) -- (1.24,1.12) -- (1.26,1.13) -- (1.28,1.14) -- (1.3,1.15) -- (1.32,1.16) -- (1.34,1.17) -- (1.36,1.18) -- (1.38,1.19) -- (1.4,1.2) -- (1.42,1.21) -- (1.44,1.22) -- (1.46,1.23) -- (1.48,1.24) -- (1.5,1.25) -- (1.52,1.26) -- (1.54,1.27) -- (1.56,1.28) -- (1.58,1.29) -- (1.6,1.3) -- (1.62,1.31) -- (1.64,1.32) -- (1.66,1.33) -- (1.68,1.34) -- (1.7,1.35) -- (1.72,1.36) -- (1.74,1.37) -- (1.76,1.38) -- (1.78,1.39) -- (1.8,1.4) -- (1.82,1.41) -- (1.84,1.42) -- (1.86,1.43) -- (1.88,1.44) -- (1.9,1.45) -- (1.92,1.46) -- (1.94,1.47) -- (1.96,1.48) -- (1.98,1.49) -- (2,1.5) -- (2.02,1.51) -- (2.04,1.52) -- (2.06,1.53) -- (2.08,1.54) -- (2.1,1.55) -- (2.12,1.56) -- (2.14,1.57) -- (2.16,1.58) -- (2.18,1.59) -- (2.2,1.6) -- (2.22,1.61) -- (2.24,1.62) -- (2.26,1.63) -- (2.28,1.64) -- (2.3,1.65) -- (2.32,1.66) -- (2.34,1.67) -- (2.36,1.68) -- (2.38,1.69) -- (2.4,1.7) -- (2.42,1.71) -- (2.44,1.72) -- (2.46,1.73) -- (2.48,1.74) -- (2.5,1.75) -- (2.52,1.76) -- (2.54,1.77) -- (2.56,1.78) -- (2.58,1.79) -- (2.6,1.8) -- (2.62,1.81) -- (2.64,1.82) -- (2.66,1.83) -- (2.68,1.84) -- (2.7,1.85) -- (2.72,1.86) -- (2.74,1.87) -- (2.76,1.88) -- (2.78,1.89) -- (2.8,1.9) -- (2.82,1.91) -- (2.84,1.92) -- (2.86,1.93) -- (2.88,1.94) -- (2.9,1.95) -- (2.92,1.96) -- (2.94,1.97) -- (2.96,1.98) -- (2.98,1.99) -- (3,2)}
\fxlcurve{vanleer}{(0,0) -- (0.02,0.039216) -- (0.04,0.076923) -- (0.06,0.113208) -- (0.08,0.148148) -- (0.1,0.181818) -- (0.12,0.214286) -- (0.14,0.245614) -- (0.16,0.275862) -- (0.18,0.305085) -- (0.2,0.333333) -- (0.22,0.360656) -- (0.24,0.387097) -- (0.26,0.412698) -- (0.28,0.4375) -- (0.3,0.461538) -- (0.32,0.484848) -- (0.34,0.507463) -- (0.36,0.529412) -- (0.38,0.550725) -- (0.4,0.571429) -- (0.42,0.591549) -- (0.44,0.611111) -- (0.46,0.630137) -- (0.48,0.648649) -- (0.5,0.666667) -- (0.52,0.684211) -- (0.54,0.701299) -- (0.56,0.717949) -- (0.58,0.734177) -- (0.6,0.75) -- (0.62,0.765432) -- (0.64,0.780488) -- (0.66,0.795181) -- (0.68,0.809524) -- (0.7,0.823529) -- (0.72,0.837209) -- (0.74,0.850575) -- (0.76,0.863636) -- (0.78,0.876404) -- (0.8,0.888889) -- (0.82,0.901099) -- (0.84,0.913043) -- (0.86,0.924731) -- (0.88,0.93617) -- (0.9,0.947368) -- (0.92,0.958333) -- (0.94,0.969072) -- (0.96,0.979592) -- (0.98,0.989899) -- (1,1) -- (1.02,1.009901) -- (1.04,1.019608) -- (1.06,1.029126) -- (1.08,1.038462) -- (1.1,1.047619) -- (1.12,1.056604) -- (1.14,1.065421) -- (1.16,1.074074) -- (1.18,1.082569) -- (1.2,1.090909) -- (1.22,1.099099) -- (1.24,1.107143) -- (1.26,1.115044) -- (1.28,1.122807) -- (1.3,1.130435) -- (1.32,1.137931) -- (1.34,1.145299) -- (1.36,1.152542) -- (1.38,1.159664) -- (1.4,1.166667) -- (1.42,1.173554) -- (1.44,1.180328) -- (1.46,1.186992) -- (1.48,1.193548) -- (1.5,1.2) -- (1.52,1.206349) -- (1.54,1.212598) -- (1.56,1.21875) -- (1.58,1.224806) -- (1.6,1.230769) -- (1.62,1.236641) -- (1.64,1.242424) -- (1.66,1.24812) -- (1.68,1.253731) -- (1.7,1.259259) -- (1.72,1.264706) -- (1.74,1.270073) -- (1.76,1.275362) -- (1.78,1.280576) -- (1.8,1.285714) -- (1.82,1.29078) -- (1.84,1.295775) -- (1.86,1.300699) -- (1.88,1.305556) -- (1.9,1.310345) -- (1.92,1.315068) -- (1.94,1.319728) -- (1.96,1.324324) -- (1.98,1.328859) -- (2,1.333333) -- (2.02,1.337748) -- (2.04,1.342105) -- (2.06,1.346405) -- (2.08,1.350649) -- (2.1,1.354839) -- (2.12,1.358974) -- (2.14,1.363057) -- (2.16,1.367089) -- (2.18,1.371069) -- (2.2,1.375) -- (2.22,1.378882) -- (2.24,1.382716) -- (2.26,1.386503) -- (2.28,1.390244) -- (2.3,1.393939) -- (2.32,1.39759) -- (2.34,1.401198) -- (2.36,1.404762) -- (2.38,1.408284) -- (2.4,1.411765) -- (2.42,1.415205) -- (2.44,1.418605) -- (2.46,1.421965) -- (2.48,1.425287) -- (2.5,1.428571) -- (2.52,1.431818) -- (2.54,1.435028) -- (2.56,1.438202) -- (2.58,1.441341) -- (2.6,1.444444) -- (2.62,1.447514) -- (2.64,1.450549) -- (2.66,1.453552) -- (2.68,1.456522) -- (2.7,1.459459) -- (2.72,1.462366) -- (2.74,1.465241) -- (2.76,1.468085) -- (2.78,1.470899) -- (2.8,1.473684) -- (2.82,1.47644) -- (2.84,1.479167) -- (2.86,1.481865) -- (2.88,1.484536) -- (2.9,1.487179) -- (2.92,1.489796) -- (2.94,1.492386) -- (2.96,1.494949) -- (2.98,1.497487) -- (3,1.5)}
\fxlcurve{superbee}{(0,0) -- (0.02,0.04) -- (0.04,0.08) -- (0.06,0.12) -- (0.08,0.16) -- (0.1,0.2) -- (0.12,0.24) -- (0.14,0.28) -- (0.16,0.32) -- (0.18,0.36) -- (0.2,0.4) -- (0.22,0.44) -- (0.24,0.48) -- (0.26,0.52) -- (0.28,0.56) -- (0.3,0.6) -- (0.32,0.64) -- (0.34,0.68) -- (0.36,0.72) -- (0.38,0.76) -- (0.4,0.8) -- (0.42,0.84) -- (0.44,0.88) -- (0.46,0.92) -- (0.48,0.96) -- (0.5,1) -- (0.52,1) -- (0.54,1) -- (0.56,1) -- (0.58,1) -- (0.6,1) -- (0.62,1) -- (0.64,1) -- (0.66,1) -- (0.68,1) -- (0.7,1) -- (0.72,1) -- (0.74,1) -- (0.76,1) -- (0.78,1) -- (0.8,1) -- (0.82,1) -- (0.84,1) -- (0.86,1) -- (0.88,1) -- (0.9,1) -- (0.92,1) -- (0.94,1) -- (0.96,1) -- (0.98,1) -- (1,1) -- (1.02,1.02) -- (1.04,1.04) -- (1.06,1.06) -- (1.08,1.08) -- (1.1,1.1) -- (1.12,1.12) -- (1.14,1.14) -- (1.16,1.16) -- (1.18,1.18) -- (1.2,1.2) -- (1.22,1.22) -- (1.24,1.24) -- (1.26,1.26) -- (1.28,1.28) -- (1.3,1.3) -- (1.32,1.32) -- (1.34,1.34) -- (1.36,1.36) -- (1.38,1.38) -- (1.4,1.4) -- (1.42,1.42) -- (1.44,1.44) -- (1.46,1.46) -- (1.48,1.48) -- (1.5,1.5) -- (1.52,1.52) -- (1.54,1.54) -- (1.56,1.56) -- (1.58,1.58) -- (1.6,1.6) -- (1.62,1.62) -- (1.64,1.64) -- (1.66,1.66) -- (1.68,1.68) -- (1.7,1.7) -- (1.72,1.72) -- (1.74,1.74) -- (1.76,1.76) -- (1.78,1.78) -- (1.8,1.8) -- (1.82,1.82) -- (1.84,1.84) -- (1.86,1.86) -- (1.88,1.88) -- (1.9,1.9) -- (1.92,1.92) -- (1.94,1.94) -- (1.96,1.96) -- (1.98,1.98) -- (2,2) -- (2.02,2) -- (2.04,2) -- (2.06,2) -- (2.08,2) -- (2.1,2) -- (2.12,2) -- (2.14,2) -- (2.16,2) -- (2.18,2) -- (2.2,2) -- (2.22,2) -- (2.24,2) -- (2.26,2) -- (2.28,2) -- (2.3,2) -- (2.32,2) -- (2.34,2) -- (2.36,2) -- (2.38,2) -- (2.4,2) -- (2.42,2) -- (2.44,2) -- (2.46,2) -- (2.48,2) -- (2.5,2) -- (2.52,2) -- (2.54,2) -- (2.56,2) -- (2.58,2) -- (2.6,2) -- (2.62,2) -- (2.64,2) -- (2.66,2) -- (2.68,2) -- (2.7,2) -- (2.72,2) -- (2.74,2) -- (2.76,2) -- (2.78,2) -- (2.8,2) -- (2.82,2) -- (2.84,2) -- (2.86,2) -- (2.88,2) -- (2.9,2) -- (2.92,2) -- (2.94,2) -- (2.96,2) -- (2.98,2) -- (3,2)}
\fxlcurve{fromm}{(0,0.5) -- (0.02,0.51) -- (0.04,0.52) -- (0.06,0.53) -- (0.08,0.54) -- (0.1,0.55) -- (0.12,0.56) -- (0.14,0.57) -- (0.16,0.58) -- (0.18,0.59) -- (0.2,0.6) -- (0.22,0.61) -- (0.24,0.62) -- (0.26,0.63) -- (0.28,0.64) -- (0.3,0.65) -- (0.32,0.66) -- (0.34,0.67) -- (0.36,0.68) -- (0.38,0.69) -- (0.4,0.7) -- (0.42,0.71) -- (0.44,0.72) -- (0.46,0.73) -- (0.48,0.74) -- (0.5,0.75) -- (0.52,0.76) -- (0.54,0.77) -- (0.56,0.78) -- (0.58,0.79) -- (0.6,0.8) -- (0.62,0.81) -- (0.64,0.82) -- (0.66,0.83) -- (0.68,0.84) -- (0.7,0.85) -- (0.72,0.86) -- (0.74,0.87) -- (0.76,0.88) -- (0.78,0.89) -- (0.8,0.9) -- (0.82,0.91) -- (0.84,0.92) -- (0.86,0.93) -- (0.88,0.94) -- (0.9,0.95) -- (0.92,0.96) -- (0.94,0.97) -- (0.96,0.98) -- (0.98,0.99) -- (1,1) -- (1.02,1.01) -- (1.04,1.02) -- (1.06,1.03) -- (1.08,1.04) -- (1.1,1.05) -- (1.12,1.06) -- (1.14,1.07) -- (1.16,1.08) -- (1.18,1.09) -- (1.2,1.1) -- (1.22,1.11) -- (1.24,1.12) -- (1.26,1.13) -- (1.28,1.14) -- (1.3,1.15) -- (1.32,1.16) -- (1.34,1.17) -- (1.36,1.18) -- (1.38,1.19) -- (1.4,1.2) -- (1.42,1.21) -- (1.44,1.22) -- (1.46,1.23) -- (1.48,1.24) -- (1.5,1.25) -- (1.52,1.26) -- (1.54,1.27) -- (1.56,1.28) -- (1.58,1.29) -- (1.6,1.3) -- (1.62,1.31) -- (1.64,1.32) -- (1.66,1.33) -- (1.68,1.34) -- (1.7,1.35) -- (1.72,1.36) -- (1.74,1.37) -- (1.76,1.38) -- (1.78,1.39) -- (1.8,1.4) -- (1.82,1.41) -- (1.84,1.42) -- (1.86,1.43) -- (1.88,1.44) -- (1.9,1.45) -- (1.92,1.46) -- (1.94,1.47) -- (1.96,1.48) -- (1.98,1.49) -- (2,1.5) -- (2.02,1.51) -- (2.04,1.52) -- (2.06,1.53) -- (2.08,1.54) -- (2.1,1.55) -- (2.12,1.56) -- (2.14,1.57) -- (2.16,1.58) -- (2.18,1.59) -- (2.2,1.6) -- (2.22,1.61) -- (2.24,1.62) -- (2.26,1.63) -- (2.28,1.64) -- (2.3,1.65) -- (2.32,1.66) -- (2.34,1.67) -- (2.36,1.68) -- (2.38,1.69) -- (2.4,1.7) -- (2.42,1.71) -- (2.44,1.72) -- (2.46,1.73) -- (2.48,1.74) -- (2.5,1.75) -- (2.52,1.76) -- (2.54,1.77) -- (2.56,1.78) -- (2.58,1.79) -- (2.6,1.8) -- (2.62,1.81) -- (2.64,1.82) -- (2.66,1.83) -- (2.68,1.84) -- (2.7,1.85) -- (2.72,1.86) -- (2.74,1.87) -- (2.76,1.88) -- (2.78,1.89) -- (2.8,1.9) -- (2.82,1.91) -- (2.84,1.92) -- (2.86,1.93) -- (2.88,1.94) -- (2.9,1.95) -- (2.92,1.96) -- (2.94,1.97) -- (2.96,1.98) -- (2.98,1.99) -- (3,2)}
\fxlstencil{4}{5.25}{3}{5.5}{4.5}{5.25}{5}{5}{2}
\fxlstencil{4}{5.25}{4.5}{5.25}{3}{5.5}{3.5}{5.25}{\tfrac{1}{2}}

    \end{tikzpicture}
    \caption{Upstream neighbour of the fine segment $f_1$ of
      \cref{fig:momentum:dual:coarse-fine} for a flux towards $+\unitvec{1}$.}
    \label{fig:momentum:limiter:stencil}
  \end{subfigure}
  \caption{Three-point face states~\cref{eq:momentum:three-point}. (a)~The
    limiters of \cref{tab:vof:limiters}, evaluated by the code at
    $\limratio{g}\in[0,3]$; Fromm and Lax--Wendroff leave the TVD region of
    Sweby~\cite{sweby1984}. (b)~Composite mesh of \cref{fig:momentum:dual}: the
    upstream neighbour $\stencildof{uu}{g}$ of the dual face
    $g=(f_1,h_1)$ is the unknown joined to $\stencildof{u}{g}=f_1$ by a dual
    face normal to $\unitvec{1}$ that lies nearest to the reflection of
    $\stencildof{d}{g}=h_1$ through $f_1$ (open circle), here the coarse face $g$
    of \cref{fig:momentum:dual:coarse-fine}; the segment $f_2$ lies on the line
    of $f_1$ and does not qualify. The spacing ratio enters the smoothness
    ratio~\cref{eq:momentum:ratio}; for the opposite flux direction the code
    selects $h_1$ with ratio~2.}
  \label{fig:momentum:limiter}
\end{figure}
